\section{Related Work}
% TODO: draft this section. No AI-generated text per team policy.
\subsection{RAG and Adaptive Retrieval}

Retrieval-Augmented Generation (RAG) combines the knowledge stored in a language model with information retrieved from an external memory. Lewis et al. introduced RAG to address the limitation that parametric language models cannot always access, update, or precisely use knowledge needed for knowledge-intensive questions. In their setup, the parametric memory is the pretrained sequence-to-sequence model, while the non-parametric memory is a dense index of Wikipedia. Retrieved passages are provided to the generator so that the answer can be conditioned on external evidence. Their results show improvements in factuality and performance on knowledge-intensive tasks.

The quality of retrieval is also important for the overall performance of a RAG system. Karpukhin et al. introduced Dense Passage Retrieval (DPR) for open-domain question answering. Unlike sparse methods such as BM25, which mainly depend on word overlap, DPR uses dense representations learned with a dual-encoder to retrieve passages based on their semantic similarity to the question. DPR improved top-20 passage retrieval accuracy over a strong BM25 baseline by 9--19 percentage points across several question answering datasets.

Although retrieval can improve the quality of an answer, it is not necessary for every query. Asai et al. note that indiscriminate retrieval can add computation and may also introduce irrelevant passages. Self-RAG addresses this by allowing the model to retrieve information on demand and to evaluate both retrieved passages and its own generation using reflection tokens. This makes retrieval more selective instead of treating it as a fixed step for every input.

Retrieval can also be performed during generation rather than only at the beginning. Jiang et al. proposed FLARE, which predicts the upcoming sentence and uses it as a query for retrieving relevant documents. When the predicted sentence contains low-confidence tokens, the retrieved information is used to regenerate the sentence. This active retrieval strategy is particularly useful for long-form and knowledge-intensive generation, where the information needed can change during generation.

Another direction is to select the retrieval strategy according to the complexity of the query. Jeong et al. proposed Adaptive-RAG, where a smaller language model predicts the complexity of an incoming query and selects between no retrieval, single-step retrieval, and iterative retrieval. Simple questions can therefore avoid unnecessary retrieval, while more complex questions can use stronger retrieval strategies. The results show that adapting the retrieval process can improve both efficiency and question answering performance.

Uncertainty can also be used as a signal for adaptive retrieval. Moskvoretskii et al. study uncertainty-based retrieval and compare it with other adaptive retrieval methods. Their analysis shows that uncertainty estimation can provide a more efficient way to decide when external information is needed, while maintaining comparable question answering performance in many settings. This supports treating retrieval as a decision that depends on the model's uncertainty rather than as a fixed operation.


\subsection{Conversational Memory}

LLMs have fixed-size context windows, which hinders their applicability in tasks like long conversation/document analysis. Conversational LLMs require previously conversed data spanning days, weeks, months, or even years to maintain consistency and personalisation. Increasing context window size becomes expensive as it increases the computational time of transformers. This issue motivated the development of persistent memory systems such as MemGPT, LongMem, MemoryBank, and LightMem. \cite{packer2023memgpt,wang2023longmem,zhong2024memorybank,fang2025lightmem}

MemGPT solved this issue by developing an OS-inspired LLM system for virtual context management. This system consists of a memory hierarchy similar to traditional OS systems, where it utilizes the concept of virtual memory paging. This system consists of main context and external context, along with the usage of a FIFO queue, recall storage, and archival storage. The FIFO queue stores the running history of user/agent messages. When the context window approaches its capacity and reaches a threshold, a memory pressure warning is created, older messages are evicted from the FIFO queue and written to recall storage. Recall storage refers to the MemGPT database, where input queries and output are written. Archival storage is the external database that stores long textual information. The main context refers to the context window that is directly accessible to the LLM, while the external context consists of recall and archival storage. The data from the external storage is retrieved using function calls when such information is required for the current conversation. This paper demonstrates how external memory can aid an LLM by providing storage beyond its context window size, which motivates our team to consider this for persistent memory. \cite{packer2023memgpt}

LongMem aids an LLM with a long-term memory bank that stores previous interactions' data. This system consists of a frozen backbone LLM, a trainable residual SideNet, and a cache memory bank. The frozen LLM encodes previous inputs, and their attention key-value pairs are stored in the cache memory. The current input in the LLM generates attention queries that are used to derive the top relevant key-value pairs from the memory bank. These key-value pairs are later fused with the current context with the aid of SideNet's memory-augmented layer. The backbone is frozen so that the representations used to create the stored memories remain consistent. The SideNet is trained separately to retrieve and read the memory, avoiding cached representations becoming stale as the backbone changes. LongMem demonstrates how relevant data based on past interactions can be fetched from external data, holding the potential to extend our system and implement persistent memory. \cite{wang2023longmem}

LLMs lack long-term memory for sustained interactions, especially when they need to remember user history, preferences, and personality. This gap motivated the development of MemoryBank, which gives LLMs long-term memory and the ability to build the user portrait from previous conversations. MemoryBank consists of three main units: Memory Storage, Memory Retrieval, and Memory Updating. Memory storage stores past conversations, summaries of events, and evolving user-personality information. MemoryBank encodes the memory pieces into vectors and retrieves memories using dense retrieval. It uses the concept of Ebbinghaus' forgetting curve, in which memory decreases over time, while recalling a memory will reinforce it and improve memory retention. This paper illustrates how conversational memory can involve retrieving past information, updating, strengthening, and forgetting data, adding additional considerations for a future persistent memory component. \cite{zhong2024memorybank}

Existing memory models pose large-scale computational and time issues, which we see in particular with the processing of redundant interaction data and performance of memory maintenance during inference. That is the issue LightMem addresses. LightMem is a lightweight and efficient memory system that tries to balance memory performance with computational efficiency. Its design is inspired by the Atkinson-Shiffrin model of human memory. Its three-stage architecture includes sensory memory, Topic-aware Short-Term Memory, and Long-Term Memory. Sensory Memory filters redundant information from raw input through lightweight compression and groups information by topic. Topic-aware Short-Term Memory groups similar/related data and stores them semantically in a mathematical space, organising the data more structurally. Long-Term Memory performs online memory updates; this is referred to as soft updates. Later, when the LLM is offline, the LLM writes the memory updates present in the soft memory, which are stored for the long term by consolidating/deduplicating/abstracting/resolving inconsistencies. This is referred to as a sleep-time update. It reduces token usage/API calls and runtime while maintaining or improving memory performance. This system provides motivation for future extensions involving online and offline memory processing and memory consolidation to reduce latency. \cite{fang2025lightmem}


\subsection{Turn-taking and Spoken Dialogue}

When humans speak, they understand when it is their turn to speak, but for conversational AIs, this is challenging. Traditionally, silence-based turn-taking models cause delays and interruptions. This motivated the engineering of turn-taking systems such as TurnGPT and VAP discussed in the proposed turn-taking systems paper. \cite{skantze2021turntaking}

The proposed turn-taking system explains how multiple linguistic, prosodic, visual, and contextual cues support continuous processing. It uses syntactic completion, pragmatic completion, prosody, pauses/fillers, gaze, head gestures, and breathing. The system can distinguish between end-of-turn detection and end-of-turn prediction. Continuous models process rolling speech and can predict upcoming completion. This paper strongly establishes why silence-based turn-taking is inefficient and can be enhanced using systems like TurnGPT and VAP. \cite{skantze2021turntaking}

TurnGPT is a Transformer-based language model that continuously processes the user's expression and predicts the probability of when the user will stop speaking after each word uttered by them. TurnGPT is aided by a GPT-2-based unidirectional Transformer which is fine-tuned on dialogue data. A Turn-Shift token is used to mark the boundary between consecutive turns occurring between the user and the LLM. TurnGPT is text-based and uses lexical information, syntactic information, dialogue history, and pragmatic information, improving the prediction of turn completion from the context of the utterance. However, TurnGPT relies on textual data and cannot process other styles of communication such as prosody and gaze. This provides a potential direction for future improvements to help our conversational LLM decide when it is its turn to speak rather than relying solely on silence-based turn-taking. \cite{ekstedt2020turngpt}

During a conversation, before the actual turn-shift, silence might occur before the user has completed their utterance, and this can be mistaken for a turn-shift. Voice Activity Projection (VAP) is a system developed to close the gap. VAP predicts future voice activity. The model predicts future voice activity of the user by processing data from the speaker's speech waveform, voice activity, and their voice-activity history. VAP can effectively determine whether the silence is speech at hold or there is supposed to be a conversation switch. It can also predict when the switch will occur, if the user will utter a short feedback, improving turn-taking conversational tasks. This provides a potential direction for using inputs beyond just textual data to make turn-taking shift prediction more reliable in future work. \cite{ekstedt2022vap}


\subsection{Decision-Theoretic Control}

Decision-theoretic approaches select actions by considering their expected outcomes and costs. Instead of treating every possible action as equally desirable, a system can compare the expected benefit of an action with the cost of taking it. This provides a general way to make decisions under uncertainty, where the preferred action depends on both the expected outcome and the information available to the system.

A related idea is selective prediction, where a model can choose not to make a prediction when its confidence is low. This creates a trade-off between the accuracy of the predictions that are made and the number of cases on which the model chooses to act. The same idea can be applied to retrieval: when the model is sufficiently confident, it may answer directly, while low confidence can provide a reason to obtain additional information.

The Value of Information (VOI) provides a decision-theoretic way to determine whether obtaining additional information is worth its cost. Information has value when it can change a decision or improve its expected outcome, but acquiring that information can also introduce a cost. For retrieval systems, this gives a natural way to view external retrieval: retrieval may improve the answer, but it also adds latency and computational cost.

Decision theory has also been applied directly to information retrieval. Cooper argued that the evaluation of information-retrieval systems should be considered in terms of expected utility, including both the value and cost of the system. This provides an early foundation for treating information retrieval as a decision problem rather than only as a fixed sequence of operations.

Similar decision-theoretic ideas have been applied to spoken dialogue. Bohus and Horvitz presented a decision-theoretic approach to turn-taking in multiparty conversation, combining uncertainty about conversational state with the costs of different outcomes. Their approach also considers uncertainty caused by computational delays in perception and response generation. This shows that conversational timing can be treated as a decision where acting too early and waiting too long can have different costs.

These works suggest that retrieval and conversational actions can be viewed through a common decision-making perspective. However, adaptive retrieval, uncertainty-based retrieval, and decision-theoretic dialogue systems have generally focused on individual decisions or a particular action. Retrieval, timing, memory use, and speaking decisions can affect one another in a conversational system. This leaves an open question of whether these actions can be selected jointly under a common utility objective.
